% Affine radix, three-digit targets, and two independent unit carries.
% Prepared for integration after the round32 certificate and proof audits.
\section{Affine radix scaling and a certificate of 95 operations}
\label{sec:affine-radix}

An affine radix saves one further operation. Its first mask forbids
one fixed bit in each coordinate digit. Three coordinates with that
bit clear represent an arbitrary nonnegative circuit value. The
larger possible digits require three-digit target windows; reset
positions make paired zero equations exact, and two different
padding carries exclude every nonunit homogeneous scaling.

Start with the system of Section~\ref{sec:linear-radix}, replacing
its two positive gaps by the combined positive gap
\begin{equation}\label{eq:affine-radix-gap}
 \ell+e+\alpha=q.
\end{equation}
This preserves the two additions used for the gaps and removes
one positive unknown and one equation. Indeed, it implies
$\ell,e<q$ and
\[
 S_2=\ell+eq\le(e+1)(q-1)\le(q-1)^2<q^2.
\]
Conversely, the canonical codes below satisfy $\ell+e<q$.
Now replace the radix and the first packed mask by
\begin{equation}\label{eq:affine-radix-change}
 H=H_0+1,\qquad B=H+b,\qquad q^2-1-b\ell,
\end{equation}
where $H_0$ is a sufficiently large fixed power of two. The
supplied fixed index remains $(V,H,T_L)$; $H_0=H-1$ is not an
additional parameter. Every Pell equation is retained.

\begin{theorem}\label{thm:best95}
Every recursively enumerable subset of the positive integers has
an admissible fixed index $(V,H,T_L)$ for a system with 34 positive
unknowns and 22 equations admitting a straight-line certificate of
95 operations: 51 multiplications and 44 additions. Here
$T_L=\psi_4(L)$ for the underlying exponent in the index
construction. The supplied parameters are $(x,V,H,T_L)$, with
three fixed index parameters besides the positive input $x$.
Fixed numerals and equality tests are free.
\end{theorem}

The following proof includes the rebuilt coefficient index and
all positive witnesses. Appendix~\ref{app:95} gives the complete
arithmetic table. The exact source and primitive checker is in
\file{../verification/}:
\begin{center}
\small\file{round32\_1980\_affine\_radix\_certificate.py}.
\end{center}

\subsection{Logical values and physical coordinates}

Compile the represented finite Diophantine system into a circuit
over nonnegative integers with actual positive input $x$,
addition, multiplication, zero and one. Signed integer witnesses
can first be represented by differences of nonnegative witnesses.
Introduce a homogeneous coordinate $\delta$, retained as one
physical coordinate. Every other logical non-input coordinate is
the sum of three distinct nonnegative physical coordinates.
Different logical coordinates have disjoint groups. The input $x$
is unsplit and has weight zero.

Every nonnegative integer $z$ has such a representation with the
$H_0$ bit clear in each entry. If that bit of $z$ is zero, take
$(z,0,0)$. If it is one, take
\begin{equation}\label{eq:affine-radix-split}
 (z-H_0,\ H_0/2,\ H_0/2).
\end{equation}
All entries are nonnegative and have the required bit clear.
They become base-$B$ digits once $B$ exceeds these finitely many
values. Their choice depends on $z,H_0$, not on how much larger
$B$ is subsequently chosen. In the reverse direction, the sum
of three decoded physical digits is a nonnegative logical value;
no further bound on that sum is needed.

Use fresh outputs and fresh square-difference copies whenever a
product would repeat a logical factor. For multiplication,
addition and zero equations use respectively
\begin{equation}\label{eq:affine-radix-gates}
 2XY-2W\delta=0,\qquad
 2(X+Y-W)\delta=0,\qquad 2X\delta=0.
\end{equation}
A copy $X'=X$ can be imposed by
\begin{equation}\label{eq:affine-radix-copy}
 X'^2-X^2=0.
\end{equation}
Nonnegativity makes this exact. Constant-one occurrences use
$\delta$ or a fresh logical copy of $\delta$. Whenever a use
in \eqref{eq:affine-radix-gates} would create $\delta\delta$,
use a fresh copy instead. Repeated operands are copied too.
Every displayed product thus has distinct logical factors and
expands into distinct physical factors. Final equalities may
use square differences or distinct-factor products with
$\delta$. Include both $F$ and $-F$ for every zero equation
$F=0$.

Introduce fresh logical coordinates $X,Y,Z,\delta',u$ with
disjoint physical groups and include both signs of
\begin{equation}\label{eq:affine-radix-normalization}
 \begin{gathered}
 X^2-x^2=0,\quad Y^2-x^2=0,\quad Z^2-x^2=0,\quad
 \delta'^2-\delta^2=0,\\
 2xX-2\delta\delta'-2u\delta=0.
 \end{gathered}
\end{equation}
Their exact zero tests will give
\begin{equation}\label{eq:affine-radix-guard}
 X=Y=Z=x,\qquad \delta'=\delta,\qquad
 x^2=\delta^2+u\delta.
\end{equation}
Consequently $\delta>0$ and $\delta\le x$. Conversely, every
positive input has the nonnegative normalization witnesses
\begin{equation}\label{eq:affine-radix-unit-witness}
 \delta=\delta'=1,\quad X=Y=Z=x,\quad u=x^2-1.
\end{equation}

Add three unpaired targets whose raw polynomial is $\delta^2$,
and put them last. The first is ordinary; the second and third
will have different padding at their preceding positions.

Expand each logical row in the physical coordinates and $x$.
Divide square coefficients by one and distinct-coordinate product
coefficients by two. Every divided coefficient lies in
$\{-1,0,1\}$. A square of a three-coordinate sum has diagonal
coefficients one and cross coefficients two. Disjoint groups and
fresh copies prevent two logical monomials from contributing to
the same physical monomial in
\eqref{eq:affine-radix-gates}--\eqref{eq:affine-radix-normalization}.
The three unit targets have divided coefficient one. These rows
and physical coordinates affect the fixed index, not the number
of unknowns of the final system.

\subsection{Support, reset positions, and the fixed index}

Let $m$ count all true non-input physical coordinates, including
the unsplit $\delta$. Assign weights
\[
 v_i=6\,3^i\quad(0\le i<m),\qquad M=6\,3^{m-1}.
\]
The input has weight zero. Quadratic monomial weights are
multiples of six, at most $2M$, and determine their unordered
coordinate pair, including squares and products with $x$.
Divide by six and read the base-three digits to see this.

Let $s$ count all targets, including both signs of every zero
equation and the three final unit targets. Define
\begin{equation}\label{eq:affine-radix-support}
 \begin{gathered}
 d_0=4M+6,\qquad
 t_j=(s+1)d_0+2M+jd_0\quad(0\le j<s),\\
 t_{\max}=t_{s-1}=2sd_0+2M,\qquad K=t_{\max}+3.
 \end{gathered}
\end{equation}
All targets are multiples of six. Write
$\tau_5=t_{s-2}$ and $\tau_7=t_{s-1}$ for the two modified
unit targets. For a target $G_j$, place divided coefficient $a$
of monomial weight $w$ at exponent $t_j-w$ in
$\mathcal D_{\rm main}(T)$. The intervals $[t_j-2M,t_j]$
are disjoint. Multiplication by the physical coordinate
polynomial squared recovers exactly $G_j$ at $t_j$.

Add the resets
\begin{equation}\label{eq:affine-radix-resets}
 \mathcal D_{\rm reset}(T)=\sum_j T^{t_j-3}.
\end{equation}
Let $\mathcal P(X)$ be the three physical indices belonging to
logical coordinate $X$, and set
\begin{equation}\label{eq:affine-radix-padding}
 \begin{aligned}
 \mathcal D_5(T)&=T^{\tau_5-1}
       +\sum_{p\in\mathcal P(X)\cup\mathcal P(Y)}
                                      T^{\tau_5-1-v_p},\\
 \mathcal D_7(T)&=T^{\tau_7-1}
       +\sum_{p\in\mathcal P(X)\cup\mathcal P(Y)\cup\mathcal P(Z)}
                                      T^{\tau_7-1-v_p},\\
 \mathcal D(T)&=\mathcal D_{\rm main}(T)
       +\mathcal D_{\rm reset}(T)+\mathcal D_5(T)+\mathcal D_7(T).
 \end{aligned}
\end{equation}
The three types have exponent residues zero, three, and five
modulo six. Each padding polynomial has distinct exponents, and
their intervals are disjoint because $\tau_7-\tau_5=d_0>M$.
All $\mathcal D$ coefficients remain in $\{-1,0,1\}$.
Every exponent is positive, greater than $M$, and below $K$.

Allow dummy physical digits at all three tested positions
$t_j,t_j+1,t_j+2$, and define
\begin{equation}\label{eq:affine-radix-codes}
 \begin{aligned}
 \ell_0(T)&=\sum_iT^{v_i}
              +\sum_j(T^{t_j}+T^{t_j+1}+T^{t_j+2}),\\
 e_0(T)&=\sum_{h=0}^{K-1}T^h+\mathcal D(T).
 \end{aligned}
\end{equation}
The digits of $\ell_0$ are zero or one, those of $e_0$ are zero,
one or two, and the unit digit of $e_0$ is one. Put
\[
 c_*=1+m+3s,\qquad D_1=\sum_h|[T^h]\mathcal D(T)|.
\]
Choose a power of two $L>3K+2$. After the complete layout is
fixed, choose a power of two $H_0$ satisfying
\begin{equation}\label{eq:affine-radix-index}
 \begin{gathered}
 H_0>\max\{2\,4^{2L+1},\ 128\max(1,D_1)c_*^2,\ 3L,\ 1024\},\\
 H=H_0+1,\qquad V=\ell_0(4)+e_0(4)4^L,\qquad T_L=\psi_4(L).
 \end{gathered}
\end{equation}
These choices precede the queried input. The fixed integer $H$
need not be a power of two.

No decoded dummy contributes to
$\mathcal D(T)\mathcal C(T)^2$ through $t_{\max}+2$.
The least $\mathcal D$ exponent is at least $t_0-2M$, the
least dummy weight is $t_0$, and
\begin{equation}\label{eq:affine-radix-dummies}
 2t_0-2M>t_{\max}+2.
\end{equation}
Thus the residue analysis below involves only true coordinates
of weight zero modulo six. This exclusion is needed because
some dummy weights are not zero modulo six.

The spacing also makes the relevant padding coefficients exact.
At $t_j-3$ the own reset contributes $x^2$. Another reset would
require a monomial weight of absolute value at least
$d_0>2M$, so cannot contribute. At $\tau_5-1$, the extra terms
give exactly $x^2+2xX+2xY$; at $\tau_7-1$, they give exactly
$x^2+2xX+2xY+2xZ$. Other targets' preceding positions receive
no extra term. Monomial uniqueness identifies these terms, the
different bands are separated by more than $4M$, and the other
parts of $\mathcal D$ have different residues modulo six.

\subsection{Preliminary bounds and the Pell conclusions}

The source definitions and coding equations are
\begin{equation}\label{eq:affine-radix-packing}
 \begin{gathered}
 \mathcal C=x+g,\quad b=x+\beta,\quad B=H+b,\quad\theta=B-4,\\
 \lambda(B-1)=q^2-1,\quad
 S_2=\ell+eq=V+t\theta,\quad\ell+e+\alpha=q,\\
 \Omega=\lambda-e>0,\quad
 \sigma=\Omega(q-\mathcal C^2)>0,\quad n=q^8,\\
 S=g+q^2(S_2+q^2\sigma),\\
 T+1=q^2(1+\theta\lambda)-b\ell+\theta\ell q^4,\\
 r=S(n^2-n)+(T+1)(n^2-1).
 \end{gathered}
\end{equation}
All supplied unknowns, including $\alpha,\beta,\Omega,\sigma$,
are positive integers. Before using any Pell conclusion or
decoded coefficient row, positivity gives
\begin{equation}\label{eq:affine-radix-preliminary}
 \ell,e<q,\quad S_2<q^2,\quad\mathcal C^2<q,\quad g<q,
 \quad0<\sigma<\lambda q<q^3.
\end{equation}
The combined gap even gives $S_2\le(q-1)^2$. Geometry gives
$B\le q^2$. Since $b\ge2$ and $H>1024$, we have $q\ge6$,
$b<B<n$, $3L\le B\le n$, and
\[
 \theta\lambda\ge B-4=H+b-4>b.
\]
Integer gaps give
\[
 0<S\le q-1+q^2(q^2-1)+q^4(q^3-1)<q^7<n.
\]
For the other packed number,
\begin{equation}\label{eq:affine-radix-packed-range}
 \begin{aligned}
 T+1&>q^2\theta\lambda-b\ell>bq^2-bq>0,\\
 T+1&<q^4(1+\theta\ell)<Bq^5\le q^7<n.
 \end{aligned}
\end{equation}
The upper bound uses $1+\theta\lambda<q^2$ and
$1+\theta\ell\le1+(B-4)(q-1)<Bq$. Consequently
\begin{equation}\label{eq:affine-radix-r-range}
 n\le n^2-1\le r<2n^3.
\end{equation}

These are the packing and growth hypotheses of the retained Pell
block. To specify the proof order, write
\[
 \begin{gathered}
 U=wn^2,\quad Y_{\rm p}=s_{\rm p}n^2,\quad
 a=Y_{\rm p}(U+1),\quad A=a+4,\\
 D_{\rm p}=A^2-1,\quad P=2UY_{\rm p}^2+1,\quad J=2r+1.
 \end{gathered}
\]
The subscripts distinguish Pell values from encoded logical
coordinates. We have $U,Y_{\rm p}\ge n^2$,
$UY_{\rm p}\ge n^4>r+1$, and $a>n^4>J$. The positive
interval and first index equation imply $c>J$. Classification
of the main and first norms, together with $P>A$, gives
preliminary main index at least $r+2$. Thus $c>AD_{\rm p}^2$,
as in Sections~\ref{sec:composed-97} and
\ref{sec:relaxed-auxiliary}.

The retained relaxed norm
\[
 (ic^2)^2=D_{\rm p}(f^2-1)
\]
has the required fundamental-unit and divisibility hypotheses.
The doubled signed-index argument gives $c=\psi_A(J)$ and
$d=\chi_A(J)$. The first-index and ratio argument gives
$Y_{\rm p}\ge U^r$ and $a>U^{r+1}$ before either exponent
relation is used. The common-witness exponent equation gives
$U=4^J$. The fixed congruence $\kappa=T_L+\Delta a$
determines the second index as $L$, and its exponent equation
gives $q=B^L$. All exponent-size inequalities follow from
\eqref{eq:affine-radix-r-range} and $3L\le B\le n\le r$
as in those proofs.

Since $n^2\mid U=4^J$, $q$ is a power of two; $q=B^L$ then
makes $B$ a power of two. From $B>H_0+1$ we get $B\ge2H_0$.
Neither $b$ nor $H$ is asserted to be a power of two, and no
Pell step requires it. The retained upper ratio and binomial
fractional-part argument gives
\begin{equation}\label{eq:affine-radix-binomial}
 n^2\mid\binom{2r}{r}.
\end{equation}
All these conclusions precede canonical decoding, internal
equations, and the proof that $\delta=1$.

\subsection{Canonical codes and the forbidden bit}

After the Pell conclusions, $b<B<q$ and $\ell<q$, so
\begin{equation}\label{eq:affine-radix-no-borrow}
 q^2-1-b\ell>0.
\end{equation}
The exact block expansion is
\[
 T=(q^2-1-b\ell)+q^2\theta\lambda+q^4\theta\ell,
\]
with widths $(2,2,4)$ in radix $q$. The binary no-carry
packing lemma, using \eqref{eq:affine-radix-packing},
\eqref{eq:affine-radix-r-range}, and
\eqref{eq:affine-radix-binomial}, gives
\begin{equation}\label{eq:affine-radix-masks}
 g\mathbin{\&}(q^2-1-b\ell)=0,\qquad
 S_2\mathbin{\&}\theta\lambda=0,\qquad
 \sigma\mathbin{\&}\theta\ell=0.
\end{equation}
The ampersand is bitwise intersection, used only to interpret
the polynomial system.

Since $q=B^L$ and $\theta=B-4$, the middle mask bounds each
base-$B$ digit of $S_2$ by three. Its digit polynomial $F(T)$
has degree less than $2L$ and $F(B)=S_2$. The source congruence
gives $F(4)\equiv V\pmod{B-4}$. Both $F(4),V<4^{2L}$,
whereas \eqref{eq:affine-radix-index} makes $B-4$ larger than
their possible difference. Hence $F(4)=V$. Uniqueness of
base-four expansion yields $S_2=\ell_0(B)+e_0(B)q$.
Both $\ell,e<q$, and both fixed codes are below $B^K<q$;
quotient and remainder uniqueness gives
\begin{equation}\label{eq:affine-radix-canonical}
 \ell=\ell_0(B),\qquad e=e_0(B).
\end{equation}
This removes the quotient alias without a quadratic digit bound.

The integer $q^2-1$ has every base-$B$ digit equal to $B-1$.
Subtracting $b\ell_0(B)$ causes no borrow. At indicator
positions the resulting digit is
\[
 B-1-b=H-1=H_0,
\]
and elsewhere it is $B-1$. Since $B,H_0$ are powers of two
with $H_0<B$, the first mask says exactly that each physical
digit of $g$ has its $H_0$ bit clear and that digits outside
the indicator support vanish. Every digit is nonnegative and
below $B$. The unit position is outside the support, so $g$
has unit digit zero. Since $x<b<B$, the unit coordinate of
$\mathcal C=x+g$ is exactly the queried input. Write
\[
 \mathcal C(T)=x+\sum_i z_iT^{v_i}
       +\sum_j(d_jT^{t_j}+d'_jT^{t_j+1}+d''_jT^{t_j+2}).
\]
Every physical digit, including the unsplit $\delta$, is below
$B$. The input is below $B$ too. Therefore
\begin{equation}\label{eq:affine-radix-coefficient-bound}
 A_{\mathcal C}=\mathcal C(1)^2<(c_*B)^2<B^3/64.
\end{equation}
The last inequality follows from the fixed bound and $B\ge2H_0$.
Every raw coefficient of $\mathcal D(T)\mathcal C(T)^2$ has
absolute value at most $A_{\mathcal C}$: distinct exponents
of $\mathcal D$ select distinct nonnegative coefficients of
$\mathcal C(T)^2$, and each multiplier has absolute value at
most one.

Also $e,\mathcal C<B^K$, hence $e\mathcal C^2<q$ because
$L>3K+2$. Through position $K-1$, the raw signed coefficients
of $\sigma$ are those of $\mathcal D(T)\mathcal C(T)^2$.
Indeed,
\[
 \sigma=(\lambda-e)q+(e-\lambda)\mathcal C^2.
\]
The first summand starts at $L$, and the coefficient polynomial
of $e-\lambda$ agrees with $\mathcal D(T)$ below $K$.
Terms starting at $K$ or $L$ cannot affect lower digits.
This identifies the actual low digits even when some raw
coefficients are negative.

\subsection{Reset padding and exact three-digit zero tests}

Use the signed carry recursion
\begin{equation}\label{eq:affine-radix-carry}
 \begin{gathered}
 h_0=0,\qquad a_j=[T^j]\mathcal D(T)\mathcal C(T)^2,\\
 h_{j+1}=\left\lfloor\frac{a_j+h_j}{B}\right\rfloor,
 \qquad\mathrm{digit}_j=a_j+h_j-Bh_{j+1}.
 \end{gathered}
\end{equation}
It computes the low base-$B$ digits of the positive integer
$\sigma$. The coefficient bound gives $|h_j|<B^2/8$ by
induction throughout the tested range: the induction step has
upper bound $B^2/64+B/8+1<B^2/8$ for $B>64$.

By \eqref{eq:affine-radix-dummies}, dummies do not contribute.
Only residue classes zero, three and five modulo six can have
nonzero raw coefficients; classes one, two and four are empty.
Two consecutive empty positions of classes one and two reduce
a carry of magnitude below $B^2/8$ to $-1$ or zero, since their
combined operation is division by $B^2$ followed by the floor.
Thus the incoming carry at every reset $t_j-3$ is in
$\{-1,0\}$.

Its raw coefficient is exactly $x^2$. Since $1\le x<B$, the
sum with the incoming carry is nonnegative and below $B^2$.
The carry to $t_j-2$ is a nonnegative integer below $B$.
That position is empty, so the carry to $t_j-1$ is exactly zero.
Except at $\tau_5,\tau_7$, the raw coefficient at $t_j-1$
is also zero. Every ordinary target therefore receives zero
incoming carry. Its next two raw coefficients are zero as well,
so its three tested digits form the residue of $G_j$ modulo
$B^3$. The third mask requires each digit to be in
$\{0,1,2,3\}$.

For a paired zero row, $|G_j|<B^3/64$. A negative nonzero
value has residue modulo $B^3$ greater than $63B^3/64$;
its top digit exceeds three, so it fails the mask. Both signs
have separate resets and zero incoming carry. They can both
pass only at zero, which passes both. Every circuit equation
and every normalization equation is now exact, without assuming
a unit value. In particular \eqref{eq:affine-radix-guard} holds.

The first unpaired unit target also receives zero carry. At
$\tau_5-1$ and $\tau_7-1$, the exact padding coefficients
and the input-copy equations give $5x^2$ and $7x^2$, with
zero incoming carry. Thus the target carries are exactly
\begin{equation}\label{eq:affine-radix-two-carries}
 h_5=\left\lfloor\frac{5x^2}{B}\right\rfloor,
 \qquad h_7=\left\lfloor\frac{7x^2}{B}\right\rfloor.
\end{equation}
Both targets still have raw value $\delta^2$ and two empty
following raw positions. Their tested digits are the residues
of $\delta^2+h_5$ and $\delta^2+h_7$ modulo $B^3$.

\subsection{Two different padding carries fix the unit}

The guard gives $0<\delta\le x<B$; the unsplit physical
digit bound also gives $\delta<B$ directly. Hence
$\delta^2<B^2$. The first unit test implies
\begin{equation}\label{eq:affine-radix-unit-digits}
 \delta^2=kB+s,\qquad k,s\in\{0,1,2,3\}.
\end{equation}
Since $4\mid B$, a square cannot have $s=2$ or $3$. If $s=1$,
then $\delta^2\le3B+3<B^2/16$ for $B\ge64$, so
$\delta<B/4$. The roots of one modulo a power of two $B\ge8$
are congruent to $1,B/2-1,B/2+1,B-1$; only the first lies in
$(0,B/4)$. Thus $\delta=1$.

For $\delta>1$ the only remaining case is
\begin{equation}\label{eq:affine-radix-large-unit}
 \delta^2=kB,\qquad1\le k\le3,\qquad x^2\ge kB.
\end{equation}
For $c=5,7$, write
\[
 h_c=m_cB+r_c,\qquad0\le r_c<B.
\]
Since $x<B$, we have $h_c<cB$. The tested integer is
$(k+m_c)B+r_c<(3+c)B<B^2$. There is no reduction modulo
$B^3$ and no overflow into its third digit. Acceptance forces
\begin{equation}\label{eq:affine-radix-carry-digits}
 0\le r_5,r_7\le3,\qquad
 0\le m_5,m_7\le3-k\le2.
\end{equation}
The two floors are related. For $z=x^2/B$, using braces for
fractional part,
\begin{equation}\label{eq:affine-radix-floor-relation}
 7h_5-5h_7=5\{7z\}-7\{5z\},\qquad
 -7<7h_5-5h_7<5.
\end{equation}
Thus its absolute value is less than seven. Substituting the
quotients and remainders and using $|7r_5-5r_7|\le21$ gives
\[
 |B(7m_5-5m_7)|<28.
\]
Since $B>32$, we obtain $7m_5=5m_7$. Its nonnegative
solutions are $(m_5,m_7)=(5h,7h)$. Bound
\eqref{eq:affine-radix-carry-digits} forces both quotients to
be zero. Then $h_5=r_5\le3$, whereas
\eqref{eq:affine-radix-large-unit} gives $h_5\ge5k\ge5$,
a contradiction. Therefore $\delta=1$.

The paired circuit equations now recover the represented system
at its actual positive input. This proves sufficiency.

\subsection{Necessity and all positive supplied witnesses}

Given a solution of the represented system, assign its logical
circuit values and fresh copies, and use
\eqref{eq:affine-radix-unit-witness} for normalization.
Split each non-input logical value other than $\delta$ into
three physical coordinates as above; retain $\delta=1$
unsplit. All physical values have their $H_0$ bit clear,
including $\delta$ since $H_0>1$. Set every dummy to zero.
Each paired target has value zero and the three unit targets
have value one.

The fixed index was chosen before this input. With the physical
assignment fixed, choose $B$ an arbitrarily large power of two.
Require $B>H+x$ and $B$ larger than every physical coordinate.
Also require every raw coefficient of
$\mathcal D(T)\mathcal C(T)^2$ to have absolute value below
$B/4$, and require $7x^2<B/4$. These conditions are possible
because all those finitely many values are fixed while $B$
grows without bound. Set
\[
 \begin{gathered}
 b=B-H,\quad\beta=b-x>0,\quad
 \ell=\ell_0(B),\quad e=e_0(B),\\
 q=B^L,\qquad n=q^8.
 \end{gathered}
\]
The sum of the two code polynomials has digits at most three,
so $\ell+e<B^K<q$. The support margin gives
$\mathcal C^2<q$ and $\lambda=(q^2-1)/(B-1)>e$.
Consequently
\[
 \begin{gathered}
 \alpha=q-\ell-e,\quad\Omega=\lambda-e,\quad
 \sigma=\Omega(q-\mathcal C^2),\\
 t=\frac{\ell+eq-V}{B-4}
 \end{gathered}
\]
are positive integers. For $t$, evaluate the nonconstant
nonnegative polynomial $\ell_0(T)+e_0(T)T^L$ at $B>4$;
its difference from its value at four is positive and divisible
by $B-4$. The encoded coordinate $\delta=1$ ensures $g>0$.

The first mask holds because the physical digits have the
specified bit clear; the second holds because the combined
code digits are at most three. For the third mask, each true
variable position lies below the least $\mathcal D$ exponent
and has zero raw coefficient and zero carry. Every paired
target receives zero carry from its reset and has value zero.
The first unit target has value one. Since $7x^2<B/4$,
both special carries vanish, so those targets also have digits
$(1,0,0)$. All masks hold. The packing lemma supplies
\eqref{eq:affine-radix-binomial} for the newly determined $r$;
the preliminary bounds give its required size estimates.

Construct every positive Pell witness for this $r$ as in
Section~\ref{sec:composed-97}. Choose $U=4^J$ and
\[
 Y_{\rm p}=\left\lfloor\frac{(U+1)^{2r}}{U^r}\right\rfloor,
 \qquad w=\frac U{n^2},\qquad
 s_{\rm p}=\frac{Y_{\rm p}}{n^2},\qquad a=Y_{\rm p}(U+1).
\]
Take the main and first Pell coordinates at $J$ and $r+1$,
their positive interval and first-index quotients, and the
second Pell coordinates at $L$. That index uses the newly
fixed $T_L=\psi_4(L)$. The retained exponent and rounding
bounds give the positive exponential quotients. They require
$B,q$ to be powers of two, as they are, but place no such
requirement on $b$ or $H$.

Choose the doubled-index auxiliary exponent last as
$m_{\rm aux}=2cJ$. Set $f=\chi_A(m_{\rm aux})$,
$i_{\rm old}=\psi_A(m_{\rm aux})/c^2>0$ and
$i=D_{\rm p}i_{\rm old}$. The doubled-index construction
supplies positive $o,j$. The relaxed auxiliary norm holds
and its signed auxiliary value is unchanged. Together with
the new positive $\beta$ and combined-bound $\alpha$, all
34 supplied unknowns are positive integers. This proves
necessity and completes the proof of Theorem~\ref{thm:best95}.

\subsection{The complete arithmetic count}

The combined-bound 96-operation schedule uses a multiplication
for $B=Hb$ and a subtraction for $b-1$. The new definition
$B=H+b$ uses one addition instead of one multiplication.
Replacing $(b-1)\ell$ by $b\ell$ preserves that multiplication
and removes the subtraction. Therefore
\[
 (52\text{ multiplications},44\text{ additions})
 \longrightarrow
 (51\text{ multiplications},44\text{ additions}),
 \qquad 51+44=95.
\]
There are 34 positive unknowns and 22 source equations. The
logical coordinates, physical triples, three-digit windows,
resets, and both padding patterns change only the fixed integer
index. No numeral construction is charged. Each subtraction
can be realized by one reversed addition equation with a free
integer intermediate.

The complete arithmetic checker compares the changed radix
and first-mask expressions with all source residuals, including
the retained acyclic corrections for geometry, the second
exponent, and the relaxed signed norm. The separate sparse
encoding regression gives finite evidence for coefficient
isolation and carries. It does not replace the universal
support, unit, Pell, or positive-witness arguments above.
The result is an upper bound, with no optimality assertion.
